\section{Novena iteración de la Fase 5: La lectura unificada}
\label{sec:F5I9}

Los datos que cargamos llegan al simulador en varias formas: el GWS como json individual, el GWB como CARPETA
con su JSON y sus imágenes, el GWI como ZIP (realmente es una
carpeta) y el GWC como FICHERO suelto. Hasta ahora cada forma
tenía su propio botón y su propio camino de entrada, cuando el
formato ya declara por sí mismo qué es: el campo \texttt{tipo} del
JSON en los cartuchos de autoría y la cabecera «GWC» en el binario.

La idea general sería cambiar a un solo botón, que presente
un panel donde arrastrar los ficheros que se van a cargar, para lo que es
necesario ejecutar dos cambios, y en este orden concreto:

\begin{itemize}
	\item Cambio del programa para unificar la forma de carga. El campo
	\texttt{tipo} del JSON en los cartuchos de autoría y la cabecera
	«GWC» en el binario indican el tipo de cartucho, así que basta con
	una única función que reciba todo lo cargado y decida por el
	CONTENIDO el tipo de cartucho y el proceso a ejecutar. El usuario
	deja así de declarar la forma en el código: la detecta el
	programa, y el botón solo la comprueba.
	
	\item Cambio en la pantalla para sustituir los botones por uno
	único, que abra un panel donde soltar los ficheros a cargar, con
	dos botones auxiliares de cargar fichero y cargar directorio.
\end{itemize}

Esta iteración aborda solo el primer punto, mejorando el código que carga los
cartuchos: un solo flujo que admite cualquiera de las tres formas y decide por
el CONTENIDO, no por el botón que se pulsó para cargarlo.

El segundo punto en este momento se deja como posible mejora general del
sistema, solo cambiaría el editor gráfico usando el nuevo proceso de lectura
de cartuchos.

\subsection{El GWS para cargar el setup}
\label{sec:F5I9GWS}
Como esta iteración unifica le lectura de los ficheros, se debe incluir al 
fichero de setup. Como complemento se aborda un pequeño cambio en los ficheros 
de setup que se definieron en la primera iteración de la fase 1 (ver apartado 
\ref{sec:Primera-iteracion}). En concreto se cambia el tipo, que era "GW", 
pasando a ser "GWS", de manera unificada con el resto de formatos.

De esta manera se puede incluir en el proceso como un tipo mas de los ficheros 
tratados, y visualmente se declara de forma mucho  más específica, pues usar 
solo "GW" podría considerarse solo el prefijo de cualquier otro tipo, al usar 
\textbf{GWS} nos aseguramos de que es un tipo único.

\subsection{Las diversas formas y su detección}
\label{sec:F5I9Formas}

El nuevo núcleo de carga  usa la función (\texttt{cargarCartucho}, en
\texttt{BaseDirector.js}), que recibe una lista de ficheros y el tipo
esperado, sea la lista de una carpeta, de un ZIP descomprimido o de
una selección de ficheros sueltos, y aplica un detector de forma en
este orden:

\begin{enumerate}
	\item Si la selección trae UN fichero \texttt{.zip}: se
	descomprime en memoria (\texttt{abrirZIP}) y el resultado vuelve a
	entrar en el detector (un ZIP es una carpeta en memoria).
	
	\item Si trae UN fichero \texttt{.gwc}: carga binaria
	(\texttt{procesarGWC}), como hasta ahora.
	
	\item En cualquier otro caso: flujo de carpeta, que ya resuelve el
	JSON prioritario por tipo (GWC $>$ GWI $>$ GWB  $>$ GWS
	\footnote{Se ha incluido el GWS para unificar completamente el sistema}) 
	y despacha el flujo automático.
\end{enumerate}

La detección es por extensión y por contenido, no por el botón: el
mismo flujo sirve para el GWB en carpeta, el GWI en ZIP o descomprimido
y el GWC suelto.

\subsection{La entrada del usuario}
\label{sec:F5I9Entrada}

Los botones actuales se renombran y se mantienen, porque una
limitación del navegador impide que un \texttt{input} elija a la vez
una carpeta y ficheros sueltos, y los tres caminos resultan claros
de manejar:

\begin{itemize}
	\item \textbf{Cargar GWS}: se cambia el nombre que antes era 
	«Cargar Setup», por unificar; carga un solo fichero y
	se lo pasa al proceso.
	
	\item \textbf{Cargar GWB}: se cambia el nombre que antes era
	«Cargar cartucho», ya que así es más claro; carga un directorio
	completo y se pasan los ficheros al proceso.
	
	\item \textbf{Cargar GWI}: se mantiene; se elige un fichero
	\texttt{.zip}, se descomprime y se pasan los ficheros al proceso.
	
	\item \textbf{Cargar GWC}: se mantiene; carga un solo fichero y
	se lo pasa al proceso.
\end{itemize}

La unificación es del FLUJO, no del botón: todos los botones llaman al
MISMO \texttt{cargarCartucho}, de modo que la forma la detecta el
código, no el botón. El proceso sería:


\begin{algorithm}[H]
	\caption{Flujo de selección y validación de carga del cartucho}
	\label{alg:carga_cartucho}
	
	\KwData{Evento de pulsación de botón en la interfaz}
	\KwResult{Cartucho cargado en el sistema o notificación de error bloqueante}
	
	\BlankLine
	\tcp{1. Interacción con el usuario}
	\textbf{Al pulsar} el botón de carga\;
		$ruta \leftarrow \text{seleccionarFicheroODirectorio}()$\;
		$ficheros \leftarrow \text{cargarFicheros}(ruta)$\;
		$tipoEsperado \leftarrow \text{obtenerTipoSegunBoton}()$\;
	\BlankLine
	\BlankLine
	\tcp{2. Llamada al proceso de validación y carga}
	\Fn{\text{CargarCartucho}($ficheros, tipoEsperado$)}{
		$tipoReal \leftarrow \text{detectarTipoPorContenido}(ficheros)$\;
		
		\eIf{$tipoReal \neq tipoEsperado$}{
			Mostrar \textbf{error bloqueante} con el motivo concreto\;
			\Return{Error (cartucho no cargado)}\;
		}{
			Ejecutar carga del cartucho\;
			Mostrar $tipoReal$ en pantalla\;
			\Return{Éxito}\;
		}
	}
\end{algorithm}


Los mensajes de error se unifican en uno solo («Cartucho NO cargado:
\dots»), con el motivo concreto detrás; desaparecen los mensajes
específicos de cada botón («ZIP GWI NO cargado», «GWC NO cargado»),
que ya no tienen sentido con un flujo único.

\subsection{Mejora futura: un solo botón con panel de carga}
\label{sec:F5I9Unificacion}

Como posible mejora futura queda anotada la sustitución de los tres
botones por UNO solo que abra un PANEL de carga: una zona donde
soltar ficheros y carpetas, con botones auxiliares de elegir fichero
y elegir carpeta.

El manejo del panel no exige ningún cambio adicional en el proceso
de lectura de cartuchos, solo tendría que llamar al mismo
\texttt{cargarCartucho} con su lista de ficheros; su coste está en
la interfaz (el panel en sí y la recursión de arrastre de carpetas,
que el navegador no resuelve por sí solo). La misma detección por
contenido permitiría además la entrada por ARRASTRE sobre la
máquina.

Al ser una mejora sobre el proceso actual, se deja de momento
pendiente, anotándose en la lista de posibles mejoras; al finalizar
el proyecto se decidirán las mejoras a incluir.

\subsection{La validación despachada por tipo}
\label{sec:F5I9Validacion}

La validación se reorganiza en tres funciones en
\texttt{BaseCartucho.js}, con la convención camelCase del proyecto:

\begin{itemize}
	\item \texttt{validarGWB} y \texttt{validarGWI}: las reglas
	específicas de cada tipo, que LLAMAN a \texttt{validarComunes} al
	principio y después añaden las suyas. Hoy no hay reglas
	específicas más allá del tipo (el contrato común las cubre), así que
	de momento son paso por el común + comprobación del tipo; la
	estructura deja el hueco preparado para cuando los formatos
	diverjan.
	
	\item \texttt{validarComunes}: el contrato de ENTRADA común a
	GWB y GWI (tipo, versión, estilo, pantalla y las tres imágenes:
	img\_s obligatoria; img\_f e img\_m opcionales pero presentes si
	se declaran). Es la \texttt{validarCartucho} actual, renombrada.
\end{itemize}

El despacho ocurre DENTRO de \texttt{leerCartucho}, una vez elegido
el JSON por prioridad de tipo: así las específicas se ejecutan
SIEMPRE, venga el cartucho de carpeta o de ZIP, y los envoltorios
\texttt{leerGWB}/\texttt{leerGWI} dejan de ser código que nadie
llama. Las funciones se exponen en el \texttt{return} del módulo
(nombre externo primero, función local después), lo que permite
probarlas de forma aislada.

\subsection{Qué cambia para el usuario}
\label{sec:F5I9Usuario}

Dos cambios visibles: los botones se
renombran de forma mas clara (\textbf{Cargar GWB}, 
\textbf{Cargar GWI} y \textbf{Cargar GWC}), y en pantalla se
muestra el tipo de cartucho cargado. 

Internamente, ningún error de
«botón equivocado» puede colar un cartucho por el camino incorrecto
(si el tipo detectado no coincide con el esperado, el cartucho no se
carga), y los mensajes de error pasan a un formato único. El anexo
de formatos (\ref{sec:F1A01GWI}) se ajusta en su redacción: el GWI se
entrega como ZIP O COMO CARPETA (mismo contenido), y la elección por
prioridad de tipos se documenta como el mecanismo que permite que
una carpeta conviva con sus tres estados.

\subsection{Problemas detectados}
\label{sec:F5I9Problemas}

Al cargar un GWI, se ha detectado que este \textbf{tenía declarada una imagen que no se incluía en el contenedor} (img\_m). 

Tras los cambios efectuados, al cargar un GWI generado por el sistema la
validación lo bloqueaba, dando un mensaje de error «no se encuentra en la carpeta:
img\_m.png», esta es la imagen opcional de la portada que se usará en un futuro selector de cartuchos. 

La cadena real del fallo:
\texttt{formatearGWI} copiaba la sección \texttt{pantalla} del GWB tal
cual (con img\_m si el GWB la declaraba), pero \texttt{guardarGWIEnDisco} 
nunca empaquetaba esa imagen en el ZIP (una decisión anterior errónea ya 
que el GWI debe incluir siempre todo lo necesario, será el GWC el que 
elimine lo que no sea imprescindible). 

El GWI resultante declaraba una imagen ausente: roto por construcción, y
el validador de entrada lo rechazaba correctamente. La primera
corrección (no declarar img\_m en el GWI) se descartó por decisión de
diseño: el GWI es el contenedor intermedio y debe ser AUTOSUFICIENTE,
llevando SIEMPRE todo lo declarado, incluso aunque resulte redundante.

La corrección final hace lo correcto: \texttt{guardarGWI} y
\texttt{guardarGWIEnDisco} reciben ahora la portada en memoria (nuevo
acceso \texttt{imagenPortada} en BaseDirector) y la empaquetan en el
ZIP con el nombre que el JSON referencia; declarada sin imagen en
memoria, error explícito. El generador del cartucho binario decidirá 
por su cuenta donde viaja la imagen para su uso en el selector de cartuchos.